This section specifies an experimental hybrid search; it does not assert
novelty of its components or superiority over published algorithms.

\subsection{Destruction and feasibility repair}
Let $Q$ be the current feasible cycle, $S(Q)$ its visited set, and
$Q^\star$ the best accepted cycle. A random destruction removes
$k=\min\{6,|S(Q)\cap R|,|S(Q)|-3\}$ optional locations, sampled uniformly
without replacement. Mandatory locations are retained. The trajectory stops
if $k=0$. A diagnostic alternative removes locations closest in input edge
cost to a random optional seed; it is evaluated separately from uniform removal.

Destruction can violate coverage and prize collection. For the remaining
visited set $S$, define
\[
 U(S)=\{w\in W:N_w\cap S=\varnothing\},\qquad
 D(S)=\max\{0,P-\sum_{i\in S}p_i\}.
\]
For an absent location $v$, the least insertion cost is
\[
 \Delta(v,Q)=\min_{(a,b)\in Q}\{c_{av}+c_{vb}-c_{ab}\},
\]
where consecutive locations include the closing edge. With
$\bar p=\max\{1,|V|^{-1}\sum_{i\in V}p_i\}$, define
\[
 g(v,S)=|\{w\in U(S):v\in N_w\}|+
                     \frac{\min\{p_v,D(S)\}}{\bar p}.
\]
Among locations with positive gain, repair minimizes
$\xi_v\Delta(v,Q)/g(v,S)$, with independent
$\xi_v\sim\mathrm{Uniform}[0.8,1.2]$. It inserts that location in a
least-cost position, updates coverage and prize, and repeats until
$U(S)=\varnothing$ and $D(S)=0$. Ties follow location and position indices.
The random multiplier also applies to negative insertion costs; no triangle
inequality is assumed. Repair is abandoned upon exhaustion of time or eligible
locations, preserving the stored feasible incumbent.

Every returned repair is checked against the original instance. Removing
only optional locations preserves mandatory visits, inserting only absent
locations preserves distinctness, and the stopping conditions enforce
coverage and prize. These properties justify feasibility of a completed
repair; they provide no guarantee of an improving objective value.

\subsection{Ordering and acceptance}
The repaired sequence is passed to LKH version 3.0.6 with an explicit
symmetric cost matrix and the visited set fixed. Each invocation has at most
$0.2$ seconds, including matrix-file construction and subprocess execution.
A validated output replaces the order only when it improves cost. Otherwise,
including after a timeout, the repaired input order remains available.
Passive counters distinguish successful exits, subprocess timeouts, and
calls without a successful exit. These measurements must be examined before
attributing poor performance to the neighborhood rather than its implementation.

For a candidate $Q'$, the best cycle is updated only by an improvement
accepted before the outer deadline. The current cycle is replaced whenever
$c(Q')\le1.01c(Q^\star)$, using the best cost after any accepted improvement.
Every twentieth destruction attempt resets the current cycle to $Q^\star$.
The current solution can therefore worsen temporarily while the accepted
best cost is nonincreasing. Destruction size, acceptance threshold, and reset
interval are frozen before a batch. This configuration has no learned policy
or adaptive operator weighting.

\subsection{Charged initialization and control}
Both total-time arms first construct a nearest-neighbor cycle through all
visitable locations and attempt beneficial feasible optional deletions under
a soft construction cap of $0.25$ seconds. Completing and validating the
cycle can exceed that soft cap; construction and setup are charged to the
outer budget, and the fallback qualifies only before that deadline. The
supplied benchmark instances admit the all-visit fallback.

The control continues the recovered MDM-GRASP under a $60$-second deadline.
The hybrid reserves a $10$-second MDM deadline and uses the remaining outer
time for destruction and repair. Both native deadlines debit the setup and
fallback time already spent. The recovered baseline uses the isolated
decoder-index and occupied-elite-count corrections, which must be disclosed
separately from algorithmic modifications. A predeadline native export seeds
the hybrid when it improves the fallback; otherwise the fallback is retained.
No expensive offline incumbent seeds this total-time comparison.

Native cleanup may return after its search deadline. A separate safety timeout
limits the process but does not extend eligibility of solutions. Return-time
overruns and the resulting loss of time for repair are retained in the records.
The budget concerns accepted solution timestamps, not a guarantee that every
operating-system process returns within exactly $60$ seconds.

\subsection{Deadline recording and component attribution}
The prospective collector records absolute \texttt{CLOCK\_MONOTONIC}
timestamps for improvements and parent start time $t_0$. An event qualifies
when its timestamp is at most $t_0+60$, even if routine cleanup finishes later.
An earlier collector instead required the entire repair routine to return
before the deadline. Its original outputs are preserved separately, and
reconstructed event bounds are treated as sensitivity data rather than
replacements for prospective measurements.

A conditional component control replays the exact saved pre-repair cycle
under its recorded remaining allowance, holding its visited set fixed and
allocating that allowance to LKH ordering with a fixed large trial cap.
This asks whether changing visits contributes beyond reordering the same
initial support. The control is a conditional replay, not a new independent
end-to-end run. It requires original scalar cost and feasibility checks,
unchanged support, source hashes, and deadline checks. This comparison must
precede any attribution of a performance gain to coverage and prize repair.
